% left_picture.tex
% \begin{tikzpicture}
%     \node[] (imageA) at (0,0){\includegraphics[width=1.1\textwidth]{monalisa/web_mona_lisa_fittest_5500_generations-200_pop}};
%     % \node[] (imageB) at (4.85,0){\includegraphics[width=0.5\textwidth]{monalisa/web_fittest_2500gen_100_pop.png}};
%     \node[below=1pt of imageA](captionA){\scriptsize Some caption.};
%     % \node[below=1pt of imageB](captionB){\scriptsize Some caption.};
%     % \draw[step=1cm,blue,thin] (0,0) grid (6,5);
% \end{tikzpicture}

% \begin{tikzpicture}
%     \node[] (imageA) at (0,0){\includegraphics[width=0.5\textwidth]{monalisa/web_fittest_2500gen_100_pop.png}};
%     \node[] (imageB) at (4.85,0){\includegraphics[width=0.5\textwidth]{monalisa/web_fittest_2500gen_100_pop.png}};
%     \node[below=1pt of imageA](captionA){\scriptsize Some caption.};
%     \node[below=1pt of imageB](captionB){\scriptsize Some caption.};
%     % \draw[step=1cm,blue,thin] (0,0) grid (6,5);
% \end{tikzpicture}
\begin{tikzpicture}[baseline=(current bounding box.north), trim left=-0.5cm,scale=0.75]
    \def\pointA{(40,50,60)}        % lab(60, 40, 50) rgb(225, 113, 56)
    \def\pointB{(25, 30,30)}        %  % lab(30, 25, 30) rgb(115, 53, 24)
    \def\pointC{(40,-50,40)}        % lab(40, 40, -50)rgb(112, 72, 177)
    % Espacio CIELAB
    \begin{axis}[
            % at={(-1cm,0cm)},
            width=12cm, height=10cm,
            % axis lines=middle,
            ylabel near ticks,
            xlabel near ticks,
            ylabel style={
                    yshift=10pt,
                    rotate=-70
                },
            xlabel style={
                    yshift=10pt,
                    rotate=-20
                },
            xlabel={$a^*$ (green–red)}, ylabel={$b^*$ (blue/yellow)}, zlabel={$L^*$ (Luminisity)},
            xmin=0, xmax=50,
            ymin=-50, ymax=50,
            zmin=0, zmax=60,
            tick label style={font=\tiny}, % Tamaño de las etiquetas
            label style={font=\normalsize,gray}, % Tamaño de los nombres de los ejes
            % xticklabels={-50,-35,-20,...,50},
            % xtick={-50,-35,-20,...,50},
            yticklabels={-50,-35,-20,...,50},
            ytick={-50,-35,-20,...,50},
            zticklabels={0,20,...,60},
            ztick={0,20,...,60},
            % enlargelimits=5.1,
            % x label style={at={(axis description cs:70.5,-0.1)},anchor=north},
            view={45}{25},
            scale only axis, % Asegura que el axis no tenga márgenes adicionales
            % axis on top % Dibuja el axis sobre otros elementos
        ]
        % Puntos de los colores
        \addplot3[only marks, color=colorTestA, mark=square*,mark options={scale=3.25}] coordinates {\pointA};
        \addplot3[only marks, color=colorTestB, mark=square*,mark options={scale=3.25}] coordinates {\pointB};
        \addplot3[only marks, color=colorTestC, mark=square*,mark options={scale=3.25}] coordinates {\pointC};
        \addplot3[only marks, color=gray, mark=*] coordinates {(40, 50, 0)}node[above right]{\scriptsize $A'$};
        \addplot3[only marks, color=gray, mark=*] coordinates {(25, 30, 0)}node[above left]{\scriptsize $B'$};
        \addplot3[only marks, color=gray, mark=*] coordinates {(40,-50,0)}node[above left]{\scriptsize $C'$};

        % Línea de distancia entre ellos
        \addplot3[dashed] coordinates {\pointA \pointB}node[below right,colorA,pos=0.55]{\small $30.7$};; % 30.7
        \addplot3[dashed] coordinates {\pointA \pointC}node[above left,colorA,pos=0.5] {\small $54.8$};; % 54.8
        % % Points dashed on xy plane
        % \addplot3[thin,dashed,gray] coordinates {(0,50,0) (40,50,0)}; % A
        % \addplot3[thin,dashed,gray] coordinates {(40,0,0) (40,50,0)}; % A
        \addplot3[thin,dashed,gray] coordinates {(40,50,0) (40,50,60)}; % A
        % \addplot3[thin,dashed,gray] coordinates {(0,30,0) (25,30,0)}; % B
        % \addplot3[thin,dashed,gray] coordinates {(25,0,0) (25,30,0)}; % B
        \addplot3[thin,dashed,gray] coordinates {(25,30,30) (25,30,0)}; % B
        % \addplot3[thin,dashed,gray] coordinates {(0,-50,0) (40,-50,0)}; % C
        % \addplot3[thin,dashed,gray] coordinates {(40,0,0) (40,-50,0)}; % C
        \addplot3[thin,dashed,gray] coordinates {(40,-50,0) (40,-50,40)}; % C
        % % \def\pointA{(40,50,60)}        % lab(60, 40, 50) rgb(225, 113, 56)
        % % \def\pointB{(25, 30,30)}        %  % lab(30, 25, 30) rgb(115, 53, 24)
        % % \def\pointC{(40,-50,40)}        % lab(40, 40, -50)rgb(112, 72, 177)
        % % \node[right,colorA] at (axis cs:45,35,50) {$\Delta E_{ab}^*=24.495$};

        % Etiquetas
        % \node[anchor=south,colorA] at (axis cs:30,50,60) {\small $A$}; % (50, 60, 40) rgb(217, 61, 54)
        % \node[anchor=north,colorA] at (axis cs:-35,-40,60) {\small $B$}; % (50,0,-40) rgb(54, 123, 187)
        % \node[anchor=north,colorA] at (axis cs:-50,40,60) {\small $C$};  % (50, -50, 40) rgb(21, 138, 44)
    \end{axis}
    % \draw[red, thick] (current bounding box.south west) rectangle (current bounding box.north east);
    \filldraw[colorTestA] (-1,9.7) rectangle ++(4mm,4mm);
    \filldraw[colorTestB] (-1,9) rectangle ++(4mm,4mm);
    \filldraw[colorTestC] (8.4,10.1) rectangle ++(4mm,4mm);
    \node[right] at (-1.35,10.4){\scriptsize CIELAB coordinates ($L^*,a^*,b^*$)};
    \node[right] at (-0.8,9.9){\scriptsize $A (60, 40, 50)$};
    \node[right] at (-0.8,9.2){\scriptsize $B (30, 25, 30)$};
    \node[right] at (8.7,10.3){\scriptsize $C (40, 40, -50)$};

\end{tikzpicture}

\begin{equation}
    \scriptsize
    \Delta E_{94}^* = \sqrt{
        \left( \frac{\Delta L^*}{k_L S_L} \right)^2 +
        \left( \frac{\Delta C^*}{k_C S_C} \right)^2 +
        \left( \frac{\Delta H^*}{k_H S_H} \right)^2
    }
    \label{eq:CIE1994}
\end{equation}

Each individual is compared with the test image comparing pixel by pixel extracting the $L^* a^* b^*$ coordinates:

\begin{center}

    \begin{tikzpicture}[scale=0.6]

        \begin{scope}[
                xshift=120,yshift=160, every node/.append style={yslant=0.5, xslant=-1},
                yslant=0.5, xslant=-1
            ]
            % Dibujar la malla de píxeles
            \foreach \x in {1} {
                    \foreach \y in {0,1,2} {
                            \fill[backgroundColor,fill opacity=0.35] (\x,\y) rectangle (\x+1,\y+1);
                        }
                }
            % Dibujar la cuadrícula con perspectiva
            \draw[step=1cm,gray,thin] (1,0) grid (2,3);
            \node[black] at(1.5,2.5){40};
            \node[black] at(1.5,1.5){40};
            \node[black] at(1.5,0.5){-50};
        \end{scope}
        \begin{scope}[
                xshift=120,yshift=70, every node/.append style={yslant=0.5, xslant=-1},
                yslant=0.5, xslant=-1
            ]
            % Dibujar la malla de píxeles
            \foreach \x in {1} {
                    \foreach \y in {0,1,2} {
                            \fill[backgroundColor,fill opacity=0.35] (\x,\y) rectangle (\x+1,\y+1);
                        }
                }
            % Dibujar la cuadrícula con perspectiva
            \draw[step=1cm,gray,thin] (1,0) grid (2,3);
            \node[black] at(1.5,2.5){30};
            \node[black] at(1.5,1.5){25};
            \node[black] at(1.5,0.5){25};
        \end{scope}
        % Definir el plano en perspectiva
        \begin{scope}[
                yshift=0, every node/.append style={yslant=0.5, xslant=-1},
                yslant=0.5, xslant=-1
            ]
            % Dibujar la malla de píxeles
            \foreach \x in {0,1,2,3,4} {
                    \foreach \y in {0,1,2,3,4} {
                            \fill[backgroundColor,fill opacity=0.5] (\x,\y) rectangle (\x+1,\y+1);
                        }
                }

            % Dibujar la silueta (ejemplo: una "L") 
            \foreach \x/\y in {1/1, 1/2, 1/3, 2/1, 3/1} {
                    \fill[colorTestB] (\x,\y) rectangle (\x+1,\y+1);
                }

            % Dibujar la cuadrícula con perspectiva
            \draw[step=1cm,gray,thin] (0,0) grid (5,5);
            \node[white] at (1.5,3.5){1};
            \node[white] at(1.5,2.5){2};
            \node[white] at(1.5,1.5){3};
            \node[white] at(2.5,1.5){4};
            \node[white] at(3.5,1.5){5};

        \end{scope}

        % Lineas entre planos
        % Dibujar líneas conectando puntos entre ambos planos
        % \draw[step=1cm,red,thin] (0,0) grid (5,5);
        % \draw[step=1cm,blue,thin] (-4,0) grid (0,5);
        \draw[colorA, thick,Triangle-] (-1,2.2) -- (-1,5);

        \begin{scope}[
                yshift=90, every node/.append style={yslant=0.5, xslant=-1},
                yslant=0.5, xslant=-1
            ]
            % Dibujar la malla de píxeles
            \foreach \x in {0,1,2,3,4} {
                    \foreach \y in {0,1,2,3,4} {
                            \fill[backgroundColor,fill opacity=0.5] (\x,\y) rectangle (\x+1,\y+1);
                        }
                }

            % Dibujar la silueta (ejemplo: una "L")
            \foreach \x/\y in {1/1, 1/2, 1/3, 2/1, 3/1} {
                    \fill[colorTestC] (\x,\y) rectangle (\x+1,\y+1);
                }

            % Dibujar la cuadrícula con perspectiva
            \draw[step=1cm,gray,thin] (0,0) grid (5,5);
            \node[white] at(1.5,3.5){1};
            \node[white] at(1.5,2.5){2};
            \node[white] at(1.5,1.5){3};
            \node[white] at(2.5,1.5){4};
            \node[white] at(3.5,1.5){5};

        \end{scope}

        \draw[Triangle-,thick,colorA](3.5,4.0)node[right]{%
            % $\begin{aligned}
            %         L^* & = 40 \\a^* & = 40 \\b^* & = -50\\
            %     \end{aligned}
            % $
        }
        to[out=180,in=90] (2.5,2.5);
        \draw[Triangle-,thick,colorA](3.6,7.2)node[right]{%
            % $\begin{aligned}
            %         L^* & = 30 \\a^* & = 25 \\b^* & = 30\\
            %     \end{aligned}
            % $
        }
        to[out=180,in=90] (2.5,5.5);

        % lab(30, 25, 30) rgb(115, 53, 24)
        % lab(40, 40, -50)rgb(112, 72, 177)

    \end{tikzpicture}
\end{center}

% La diferencia de color \(\Delta E_{94}^*\) entre dos colores se calcula mediante la siguiente ecuación:
